\section{Chained immediate transfer}
\label{dist:sec:chained}

The second probe couples cells together in the most austere way available.
Take $M$ host cells, all empty but one, which holds a single particle. Each
runs the $\mu=0$ birth--catastrophe process of the previous sections, and on
rupture the released load does not enter an extracellular pool: it passes
immediately and wholesale into a fresh cell, which begins the process again
with that load as its founding census, and so on until every cell has
ruptured. There is no extracellular phase, no clearance and no re-infection,
only a chain of bursts; and since $\mu=0$, all $M$ ruptures occur almost
surely.

Serial rupture of this kind is not a fiction:
\fwd{path}{\ChPath} gives the
biological account, with its Brownian-drift and macrophage-motility caveats.
Here it is a tractable laboratory for multi-cell burst statistics,
interpolating between the single-cell theory above and the populations of \ChPop.

Let $r_i$ be the load at the $i$th rupture, $t_i$ its time, and
$T_i=t_i-t_{i-1}$ the $i$th inter-rupture interval, with $t_0=0$. Write
\begin{equation}
s=\frac{\delta}{\lambda+\delta},
\qquad
\rho=\frac{\lambda}{\lambda+\delta}=1-s.
\label{dist:eq:srho}
\end{equation}

\subsection{The key observation: event types ignore the load}
\label{dist:sec:key-observation}

When the load is $n$, the next event is a birth with probability
$\lambda n/(\lambda n+\delta n)=\rho$ and a catastrophe with probability $s$,
independent of $n$; only the holding time depends on the load, being
$\mathrm{Exp}\lrb{(\lambda+\delta)n}$. The sequence of event \emph{types} is
therefore independent coin flips, and the number of births $G$ before the
catastrophe in a cell founded with any load is
\begin{equation}
G\sim\mathrm{Geom}_0(s),
\qquad
\pr{G=g}=s\,\rho^g,\qquad g=0,1,2,\ldots,
\label{dist:eq:geomevents}
\end{equation}
the number of failures before the first success. Each cell in the chain
contributes an \emph{independent} copy $G_i$, since the process restarts
afresh at every transfer, and the rupture size of cell $i$ is its founding
load plus its own births. The first cell being founded by a single particle,
\begin{equation}
r_k = 1 + G_1 + \cdots + G_k,
\label{dist:eq:rkdecomp}
\end{equation}
with $G_1,\ldots,G_k$ independent and identically distributed. Everything
about the rupture sizes follows from this decomposition, and something has
quietly happened in reaching it: the load, the whole subject of
\cref{dist:sec:geometric}, has dropped out of the argument.

\subsection{Rupture-size laws}
\label{dist:sec:rupture-sizes}

A sum of $k$ independent geometric variables is negative binomial, so
\cref{dist:eq:rkdecomp} gives immediately
\begin{equation}
\pr{r_k=n}=\mathcal{A}^{(n)}_{k}\,s^{k}\,\rho^{\,n-1},
\qquad
\mathcal{A}^{(n)}_{k}=\binom{n+k-2}{k-1},
\quad n\ge1.
\label{dist:eq:rklaw}
\end{equation}
The first few cases show the pattern:
\begin{equation}
\pr{r_1=n}=s\rho^{n-1},
\label{dist:eq:r1}
\end{equation}
geometric, matching the one-cell law of \cref{dist:eq:burstlaw} at $\mu=0$,
and then
\begin{equation}
\pr{r_2=n}=n\,s^2\rho^{n-1},
\qquad
\pr{r_3=n}=\tfrac{n(n+1)}{2}s^3\rho^{n-1},
\qquad
\pr{r_4=n}=\tfrac{n(n+1)(n+2)}{6}s^4\rho^{n-1}.
\label{dist:eq:r234}
\end{equation}
Equivalently the coefficients satisfy
\begin{equation}
\mathcal{A}^{(n)}_{k}=\sum_{i=1}^{n}\mathcal{A}^{(i)}_{k-1},
\qquad
\mathcal{A}^{(i)}_{1}=1,
\label{dist:eq:arec}
\end{equation}
which is \cref{dist:eq:rklaw} conditioned on the value of $r_{k-1}$:
$\pr{r_k=n}=\sum_i\pr{r_{k-1}=i}\pr{G_k=n-i}$. The closed form then follows
by induction from the hockey-stick identity.

\begin{remark}[The coefficient]
\label{dist:rem:coefficient}
The coefficient in \cref{dist:eq:rklaw} is easy to miscount. The superficially
similar expression $\frac{1}{k!}\frac{(n+k-1)!}{(n-1)!}=\binom{n+k-1}{k}$ is
off by one shift, and is already contradicted by the $r_2$ and $r_3$ cases
displayed above, which \cref{dist:eq:rklaw} reproduces exactly. The correct
coefficient is $\binom{n+k-2}{k-1}$, and it has a direct reading: it counts
the ways of distributing the $n-1$ births accumulated along the chain among
its $k$ cells, each cell receiving any non-negative number --- the
stars-and-bars count for \cref{dist:eq:rkdecomp}.
\end{remark}

Further consequences, all verified numerically in
\cref{dist:app:verification}:
\begin{align}
\E{r_k}&=1+k\,\frac{\rho}{s}=1+k\,\frac{\lambda}{\delta},
&
\var{r_k}&=k\,\frac{\rho}{s^2},
\label{dist:eq:rkmoments}\\[0.3em]
\E{z^{r_k}}&=z\Bigl(\frac{s}{1-\rho z}\Bigr)^k,
&
\sum_{n\ge1}\pr{r_k=n}&=1 .
\label{dist:eq:rkgf}
\end{align}
The mean rupture size grows linearly along the chain: each cell adds one
expected burst's worth of births, $\lambda/\delta$, to what it inherited.
\Cref{dist:fig:chained} shows the laws against exact simulation, with the
inter-rupture intervals, which shorten as the inherited load grows.

\subsection{Rupture times}
\label{dist:sec:rupture-times}

The first rupture time inherits the single-cell flux. With $m$ founders the
no-rupture probability is $I(t)^m$ by \cref{dist:eq:ihatk}, since $D\equiv0$
at $\mu=0$, so
\begin{equation}
\pr{t_1>t}=I(t)^m,
\qquad
f_{t_1}^{(m)}(t)=-m\,I(t)^{m-1}I'(t)=m\,I(t)^{m-1}\delta J(t).
\label{dist:eq:t1m}
\end{equation}
For $m=1$ this is the burst-time density $\varphi=\delta J$ of
\cref{dist:eq:phi}, proper here because $b=0$. When $\mu>0$ the same
reasoning must distinguish the two fixation routes: the \emph{burst-only}
time has distribution function $(1-I(t))/(1-b)$ and density
$\delta J(t)/(1-b)$, while the \emph{general-fixation} time --- rupture or
internal extinction --- has density
\begin{equation}
\delta J(t)+\mu\,p_1(t),
\label{dist:eq:genfix}
\end{equation}
the extra term being the flux through $1\to0$: only a population of size
exactly one can die out in a single death event.

The intervals, their Laplace transforms and the joint structure of the second
rupture time are collected in \cref{dist:app:chained-transforms}, beside the
hypergeometric machinery they need.

\subsection{Where the structure breaks}
\label{dist:sec:chained-scope}

That extra term in \cref{dist:eq:genfix} is the whole difficulty. The
$\mu=0$ chain rests on \cref{dist:eq:geomevents}, and the event-type
decomposition works precisely because a cell then has one way to stop: it
ruptures. Once $\mu>0$ a lineage can vanish internally, event types stop being
independent of the load, and \cref{dist:eq:rkdecomp} has no analogue. The
joint laws of $(t_k,r_k)$ for general $\mu>0$ therefore remain open;
\cref{dist:app:extract} sets out the available route and
\cref{dist:sec:open-problems} states what is missing. For $\mu=0$, by
contrast, \cref{dist:eq:rkdecomp} extends the size laws to all $M$ cells at
once, nothing in it depending on where a cell sits in the chain.

This contrast is why both probes were run. $G_k=G^k$ needs only independence
between lineages, and multiplicity supplies it; the chain's decomposition
needs the stopping mechanism to be unique, and internal extinction destroys
that.

\begin{figure}[tbp]
\centering
\includegraphics[width=0.92\textwidth]{fig14_chained_transfer/fig14}
\caption{Chained immediate transfer with $(\lambda,\delta)=(1,1)$ and
$s=\rho=1/2$. (a) Simulated rupture-size probabilities (bars) from $10^6$
chain realisations with seed 42, with the exact negative-binomial laws
\cref{dist:eq:rklaw} (curves) for $k=1,\ldots,4$. (b) The analytic mean
interval $\E{T(k)}$ of \cref{dist:eq:meant} at fixed founding load $k$, and
the simulated mean interval at chain position $k$; both decrease over
$k=1,\ldots,6$, the latter faster because its founding loads are random and
growing.}
\label{dist:fig:chained}
\end{figure}
